\subsection{Protocol and Artifacts}
We report completed, saved runs rather than projected performance. Four deterministic synthetic maps represent clustered post-disaster tasks, a curved wildfire front, a branching pipeline corridor, and a uniform agriculture field. All have three agents, axis limits of $3\,$m/s and $1.5\,$m/s$^2$, and a $0.5\,$s handshake. The results below use the archived threefold-flight-budget variants of these maps. Their task positions, rewards, motion limits, and communication settings are unchanged from the source maps; only each agent's flight-time budget is tripled. The maps contain 64, 55, 48, and 180 tasks, respectively. Their geometry is stylized and contains no obstacles.

Each map receives one MOVNS run with seed 42, a 115-second optimizer allowance, a 120-second process timeout, and limits of one million candidate evaluations and one million outer iterations. The first limit reached stops the run. The archive epsilon widths are one reward unit, $0.1\,$s mean time, and $1\,$s positive communication violation. Free waypoints are capped at three per agent. Two maps run concurrently. Full configuration, source hashes, trajectories, archives, logs, and plots are saved in the repository's \texttt{results/step4/scenarios\_search\_guidance\_triple\_budget\_20261012\_4maps} folder. An independent reconstruction checker validated all four saved epsilon fronts.

\subsection{Observed Fronts}
Table~\ref{tab:scenarios} reports the highest-reward member of each zero-violation front, breaking reward ties by lower mean completion time. The last reward column is the highest-reward \emph{saved} motion- and budget-valid candidate when the communication bound is ignored; it is a diagnostic comparator, not an optimum or a formal upper bound. The values are taken from the saved result and representative-path metadata.

\begin{table}[t]
\centering
\caption{Single-seed, 115-second searches on threefold-flight-budget maps. $Q_f$ is the best saved $T_c$-feasible reward and $Q_b$ the best saved budget-only reward.}
\label{tab:scenarios}
\resizebox{\columnwidth}{!}{%
\begin{tabular}{lrrrrrrrr}
\toprule
Map & $B_m$ & $R$ & $T_c$ & $Q_f$ & $\bar c$ & $\Delta$ & $Q_b$ & Iter.\\
\midrule
Post-disaster & 360 & 12 & 12 & 186 & 65.49 & 11.79 & 1035 & 2\\
Wildfire & 450 & 14 & 15 & 339 & 70.96 & 0.50 & 1348 & 2\\
Pipeline & 480 & 15 & 20 & 118 & 65.00 & 0.50 & 579 & 3\\
Agriculture & 465 & 18 & 25 & 30 & 44.92 & 0.50 & 367 & 1\\
\bottomrule
\end{tabular}}
\end{table}

The saved zero-violation fronts contain 16, 25, 11, and 8 epsilon representatives, respectively. Their best-reward solutions have no FREE visit or planned meeting. The observed $Q_f$--$Q_b$ gaps therefore cannot be attributed solely to the communication constraint: they also reflect incomplete search. In particular, agriculture completed just one outer iteration. The number of evaluated candidates was 1328, 1378, 1423, and 1084, respectively. These counts and the objective-evolution plots show that a short wall-time cap spends substantial effort on initialization and early routing families before communication neighborhoods receive comparable exploration.

\begin{figure}[t]
\centering
\includegraphics[width=\columnwidth]{../results/step4/scenarios_search_guidance_triple_budget_20261012_4maps/02_wildfire_front/paths-max-reward.png}
\caption{Highest-reward $T_c$-feasible wildfire route from the saved threefold-budget run. Each agent has its own speed-gradient color; collected tasks use distinct markers. The co-located pattern illustrates the current search's tendency to retain convoy-like solutions.}
\label{fig:wildfire_route}
\end{figure}

A separate, deliberately small two-agent relay probe tests whether a free waypoint can improve a feasible route. Its saved best feasible solution collects reward 146 with mean completion time $55.13\,$s, peak staleness $9.11\,$s below $T_c=12\,$s, and a FREE visit at $(67.4,0)$. A saved same-reward feasible candidate without a FREE visit has mean time $57.65\,$s. This is a within-run example of benefit, not a controlled ablation or a general performance estimate.

The budget multiplier is not a monotonicity claim about the heuristic. A larger physical budget enlarges the feasible set, but a fixed-time stochastic search need not return a better incumbent. The archived original-versus-tripled comparison summary illustrates this: post-disaster's best saved feasible reward fell from 311 to 186 when the larger budget changed the search trajectory. We do not infer a causal improvement or deterioration from one seed. Matched multi-seed comparisons, longer search budgets, communication-neighborhood ablations, and a kinematically matched TOP-CC baseline remain future experimental work.
